% Standalone demo-introduction slides for 2026-05-27.

\begin{frame}{Background: RCA and Long-Context Evidence}
  \begin{itemize}
    \item \textbf{RCA}: root-cause analysis identifies the failing service,
      fault type, supporting evidence, confidence, severity, and recommended fix.
    \item \textbf{Evidence bundle}: an RCA case can include incident symptoms,
      service topology, logs, traces, fault reports, and historical context.
    \item \textbf{Long-context challenge}: useful evidence may be spread across
      many noisy chunks, while direct prompting increases context and KV cost.
    \item \textbf{Compressed-memory RCA}: the model should answer from a compact
      memory state while preserving root-cause quality and evidence grounding.
  \end{itemize}
\end{frame}

\begin{frame}{Significance: What This Proves}
  \begin{itemize}
    \item \textbf{For the RCA direction}: validate whether open-source LLMs can
      learn RCA structure from public RCA datasets.
    \item \textbf{For product use}: expose model behavior through a chat probe,
      not only through aggregate metric tables.
    \item \textbf{For the next research step}: decide whether long-context
      compression is needed beyond direct SFT.
    \item \textbf{For release safety}: separate public RCA data and private
      Nokia DTS probes so public artifacts remain clean.
  \end{itemize}
\end{frame}

\begin{frame}{What We Built: RCA SFT Matrix}
  \begin{itemize}
    \item Built a matrix over two base families: Qwen2.5-7B-Instruct and Qwen3-8B.
    \item Compared base, curriculum, joint-mix, and anti-forget variants.
    \item Evaluated on four public RCA datasets:
      \begin{itemize}
        \item Nezha;
        \item OpenRCA-500;
        \item RCAEval;
        \item lincyaw/rca.
      \end{itemize}
    \item Added five general benchmarks to measure whether RCA SFT damages
      non-RCA capability.
  \end{itemize}
\end{frame}

\begin{frame}{What We Built: RCA Chat Probe}
  \begin{itemize}
    \item Built `rca-chat`, a Gradio probe for loading matrix checkpoints.
    \item The probe supports base, curriculum, and mix checkpoints from the RCA
      matrix.
    \item It can sample public RCA cases and optionally load private Liang DTS
      cases for internal-only qualitative inspection.
    \item The purpose is to inspect behavior that tables hide:
      \begin{itemize}
        \item JSON format stability;
        \item evidence citation;
        \item service-level root-cause prediction;
        \item qualitative answer usefulness.
      \end{itemize}
  \end{itemize}
\end{frame}

\begin{frame}{What We Learned From the Demo}
  \begin{itemize}
    \item SFT can teach useful RCA signals: service identification,
      fault/severity fields, recommendations, and structured JSON behavior.
    \item Direct SFT is not always robust:
      \begin{itemize}
        \item some checkpoints learn RCA but lose output format;
        \item anti-forget v2 improved surface JSON but lost evidence grounding;
        \item general capability can move when training only on RCA data.
      \end{itemize}
    \item Conclusion: the RCA demo supports SFT as an initialization step, but
      not as the final long-context RCA solution.
  \end{itemize}
\end{frame}

\begin{frame}{Next Demo: Long-Context RCA Compression}
  \begin{itemize}
    \item Next, use one public RCA case to compare four input strategies:
      \begin{enumerate}
        \item full context;
        \item summary;
        \item retrieval;
        \item compressed-memory RCA.
      \end{enumerate}
    \item Report the same axes across all four strategies:
      \begin{itemize}
        \item RCA answer quality;
        \item evidence grounding;
        \item output format stability;
        \item context / KV cost reduction.
      \end{itemize}
    \item Decision gate: continue toward wrapper memory only if it beats a cheap
      baseline on at least one meaningful axis.
  \end{itemize}
\end{frame}
